% ID: SEL-D0060-OB-E01
% ESTADO: cotejada. Segunda pasada visual y borrador compilado revisados.
% Fuente: Selectividad/sel_2011_matematicas_aplicadas/223A-EXAMEN 3.pdf
% SHA-256: 5998699e0b2992b988a8e0ed0d2e3581dd8fac52d6f9802587ebd40fee5ef608
% Página: 2; opción B; ejercicio original 1.

Se considera el recinto $R$ del plano determinado por las siguientes inecuaciones:
\[
13x+8y\leq600;\quad3(x-2)\geq2(y-3);\quad x-4y\leq0.
\]
\begin{enumerate}[label=\alph*)]
\item \textbf{(1.75 puntos)} Represente gráficamente el recinto $R$ y calcule sus vértices.
\item \textbf{(0.75 puntos)} Calcule el valor máximo en dicho recinto de la función $F(x,y)=65x+40y$, indicando dónde se alcanza.
\end{enumerate}
